\documentclass[10pt,a4paper]{amsart}
\usepackage[margin=19mm]{geometry}
\usepackage{amsmath,amssymb,mathtools}
\usepackage[scr=rsfs,cal=euler]{mathalfa}
\usepackage{xfrac}
\usepackage[unicode,hidelinks,bookmarks=false]{hyperref}
\newcommand{\ie}{i.e.\ }
\newcommand{\cf}{cf.\ }
\newcommand{\resp}{respectively\ }
\newcommand{\Subsection}{\subsection}
\providecommand{\nobreakdash}{\nobreak\hbox{-}}
\setlength{\emergencystretch}{2em}
\multlinegap0pt
\usepackage{bm}
\usepackage{bbm}
\usepackage{dsfont}
\usepackage{xspace}
\usepackage{cancel}
\usepackage{mathtools}
\usepackage{amsthm}
\makeatletter
\@ifundefined{theorem}{\newtheorem{theorem}{Theorem}}{}
\@ifundefined{defn}{\newtheorem{defn}[theorem]{Definition}}{}
\@ifundefined{lemma}{\newtheorem{lemma}[theorem]{Lemma}}{}
\makeatother
\title[Upper chromatic number of generalized quadrangles]{Upper chromatic number of generalized quadrangles}
\author{Gabriela Araujo-Pardo, Juli\'an Fres\'an, Linda Lesniak, Mika Olsen}
\date{Source: arXiv:2607.08887v1 (CC BY 4.0)}
\begin{document}
\maketitle

\noindent\emph{Source and licence.} Upper chromatic number of generalized quadrangles. Version arXiv:2607.08887v1 (CC BY 4.0), distributed under the Creative Commons Attribution 4.0 International licence, as recorded in the arXiv licence metadata for this exact version and shown on the abstract page https://arxiv.org/abs/2607.08887v1. Full rights notice: licenses/arxiv-2607.08887-CC-BY-4.0.txt.\par\smallskip

\noindent\textit{Editorial scope.} Counted unit: the Lemma whose source label is \texttt{Lem:CotaB\_q-1} in the article (source0.tex lines 392--394, source label \texttt{Lem:CotaB\_q-1}), delivered together with the article's own complete original proof of it (source0.tex lines 396--424, one \texttt{proof} environment of 29 lines). No mathematics is written, repaired or re-derived: statement and proof are the article's own text line for line. Required context retained uncounted: Lem:CotaB (lines 352--359); def:projection (lines 377--379). Every definition, display and label of those lines is retained unchanged. Omitted from this package and not redistributed: 1--351, 360--376, 380--391, 395--395, 425--667. Nothing inside a retained excerpt is omitted or altered: every retained body is delivered exactly as the article's lines are written, so no omission receipt of any kind accompanies this package. Citations: the retained excerpts carry no citation command at all, so no bibliography entry of the article is transcribed and no reference is redistributed. Reference adaptation: none was needed and none was made; every reference inside the delivered unit resolves inside it, so no unresolved reference or citation is produced. Printed slips kept literal: no printed word, symbol, hypothesis or number of the counted statement or of its proof was altered. Layout and class: the article's own class is \texttt{article} with packages including inputenc, graphicx, amssymb, color, amsmath, amsthm, epic, tikz, url, soul, cancel, ulem, xcolor, amscd; the wrapper is the lane's pinned \texttt{amsart} editorial wrapper at 10pt on A4, which restates the theorem-like environments the retained text needs and adds the article's own macro definitions that the retained text needs (none), with extra wrapper packages (bm, bbm, dsfont, xspace, cancel, mathtools). The delivered numbering is the unit's own. Assets: no retained span contains a figure environment, a graphics-inclusion command, a PDF-page inclusion, a table or a TikZ picture; no image file is copied into this package, the package declares no asset dependency and no image file is redistributed with it. Licence: version arXiv:2607.08887v1 of this article is distributed under the Creative Commons Attribution 4.0 International licence, as recorded in the arXiv licence metadata for this exact version and shown on the abstract page https://arxiv.org/abs/2607.08887v1. A full rights notice is delivered at licenses/arxiv-2607.08887-CC-BY-4.0.txt. Characters: every retained excerpt is pure ASCII, so the delivered engine, XeLaTeX with its default Latin Modern text font, reports no missing character.
\begin{lemma}\label{Lem:CotaB}
    Let $B$ be a 2-blocking set. If $G=(B,\mathcal{L})$ and $|B|=2(q^2+1)$ then:
    \begin{enumerate}
    \item $d_G(p)=q+1$ for every vertex $p\in B$ and $d_G(\ell)=2$ for every vertex $\ell \in \mathcal{L}$,
        \item $E^*=\emptyset$, and 
    \item For every connected component $G_s=(B_s,\mathcal{L}_s)$, $B_s$ has at least $2q+2$ points.
    \end{enumerate}
\end{lemma}   

\begin{defn}\label{def:projection}
Let $G$ be a bipartite graph with partite sets $U$ and $V$. The \emph{$U$-partite projection of $G$}, denoted as $G_U$, is a graph whose vertex set is $U$ and two distinct vertices $u_1, u_2 \in U$ are adjacent in $G_U$ if and only if $N(u_1)\cap N(u_2)\neq\emptyset$.     
\end{defn}

\begin{lemma}\label{Lem:CotaB_q-1}
  If $B$ is a 2-blocking set of order $2(q^2+1)$ such that $G=(B,\mathcal{L})$ has $q-1$ connected components, then $|B_i|\ge2q+4$ for each component $G_i=(B_i,\mathcal{L}_i)$.
\end{lemma}

\begin{proof}
     By Lemma \ref{Lem:CotaB}, $B_i$ has at least $2q+2$ points and $E^*=\emptyset$.
     
     If $G$ has $q-1$ components, for a contradiction, assume that there exists a component $G_i$ with $|B_i| = 2q + 2$. 
    %
    Since every point in $B_i$ has degree $q+1$ in the incidence graph, and every line in $L_i$ is covered by $B_i$ (it has degree 2 in $G_i$), counting the edges between $B_i$ and $\mathcal{L}_i$, we obtain $2|L_i| = |B_i|(q+1) = (2q+2)(q+1) = 2(q+1)^2$. Thus, $|L_i| = (q+1)^2 = q^2 + 2q + 1$.
    
Let $R=\mathcal{P}-B_i$ and consider the bipartite graph $H = (R, \mathcal{L}_i)$ induced by the set of points $R$ and the lines in $\mathcal{L}_i$. 
Each line $\ell \in \mathcal{L}_i$ is covering by $B_i$ (that is $N_{G_i}(\ell)=2$). Since every line in $\mathcal{Q}(4,q)$ has $q+1$ points, the remaining $q-1$ points must be in $R$. Therefore, $d_H(\ell) = q - 1$ for every line $\ell \in \mathcal{L}_i$. Then we have that:  
\begin{equation}\label{eqL}
 \sum_{l \in L_i} d_H(l) = |L_i|(q - 1) = (q^2 + 2q + 1)(q - 1) = q^3 + q^2 - q - 1.
\end{equation}
Consider the bipartite projection $G_{B_i}$ (Definition \ref{def:projection}). The incidence graph of $\mathcal{Q}(4,q)$ has girth 8, so $G_{i_{B_i}}$ has no 3-cycles. The maximum number of edges in a triangle-free graph of order $|B_i|=2q+2$ is equal to $(q+1)^2$, and it is obtained by the complete bipartite graph $K_{q+1,q+1}$. Thus, $G_{i_{B_i}}\cong K_{q+1,q+1}$ and the distance between any two lines in $G_i$ is at most 4. 

Since $\mathcal{Q}(4,q)$ has girth 8, the degree in $H$ of any point $p' \in R$ is at most 1 (i.e., $d_H(p') \le 1$). This implies that the sum of the degrees of the points in $R$ is bounded by the total number of points in $R$:
\begin{equation}\label{eqP}
    \sum_{p' \in R} d_H(p') \le |R|  = q^3 - q^2 + q - 1. 
\end{equation}
Since $H$ is a bipartite graph, the sum of the degrees of the vertices in both partitions must be equal:
$\sum_{p' \in R} d_H(p') = \sum_{\ell \in \mathcal{L}_i} d_H(\ell)$. But then Equations \ref{eqL} and \ref{eqP} lead to a contradiction.

Analogously, if $|B_i|=2q+3$, since $G_i$ is bipartite,  $$\sum_{\ell\in\mathcal{L}_i} d_{G_i}(\ell)= \sum_{p\in B_i} d_{G_i}(p) \mbox{, and } |\mathcal{L}_i|=\sum_{p\in B_i} d_{G_i}(p)/2.$$ 
It follows that $|\mathcal{L}_i|=(2q+3)(q+1)/2$.
As previously, since $\mathcal{Q}(4,q)$ has girth 8, the degree in $H$ of any point $p' \in R$ is at most 1 (i.e., $d_H(p') \le 1$),
$$ \sum_{p' \in R} d_H(p') \le |R|  = q^3 - q^2 + q - 1,\text{ and }
\sum_{l \in L_i} d_H(l) = (q - 1)(2q+3)(q+1)/2.$$
This contradicts $ \sum_{p' \in R} d_H(p') = \sum_{\ell \in \mathcal{L}_i} d_H(\ell) $.
Thus, if $q\ge3$, every $B_i$ has order at least $2q+4$.
\end{proof}

\end{document}
